\documentclass[12pt,a4paper]{article}

\usepackage[T2A]{fontenc}
\usepackage[utf8]{inputenc}
\usepackage[russian]{babel}
\usepackage{amsmath}
\usepackage{booktabs}
\usepackage{geometry}
\usepackage{pgfplots}
\usepackage{float}
\usepackage{array}

\geometry{margin=2cm}
\pgfplotsset{compat=1.18}

\begin{document}
\begin{center}
    {\LARGE Лабораторная работа 4}
\end{center}
\begin{flushright}
\textbf{Работу выполнили:} \\
Шубин Егор\\
Санатин Дмитрий\\
\textbf{Преподаватель:} \\
Танченко Юлия Валерьевна
\end{flushright}

\section*{Условие}
Дана выборка пар $(x_i,t_i)$, $i=1,\dots,n$, где $n=20$. Требуется методом
наименьших квадратов найти коэффициенты регрессионной модели
\[
    f(x)=w_0+w_1x+w_2x^2,
\]
а также вычислить среднеквадратическое отклонение
\[
    RMSE=\sqrt{\frac{1}{n}\sum_{i=1}^{n}\left(t_i-f(x_i)\right)^2}.
\]
На графике необходимо изобразить исходные точки выборки $(x_i,t_i)$ и график
полученной регрессии.

Несмотря на то, что в условии используется формулировка ``линейная
регрессия'', модель является квадратичной по переменной $x$. При этом она
остается линейной по неизвестным параметрам $w_0,w_1,w_2$, поэтому коэффициенты
можно находить стандартным методом наименьших квадратов через систему
нормальных уравнений.

\section*{Исходные данные}
\begin{center}
\begin{tabular}{ccccccccccc}
\toprule
$i$ & 1 & 2 & 3 & 4 & 5 & 6 & 7 & 8 & 9 & 10 \\
\midrule
$x_i$ & 1.0 & 1.1 & 1.2 & 1.3 & 1.4 & 1.5 & 1.6 & 1.7 & 1.8 & 1.9 \\
$t_i$ & 2.0 & 1.9 & 1.7 & 1.6 & 1.3 & 1.1 & 1.3 & 1.2 & 1.1 & 1.0 \\
\bottomrule
\end{tabular}
\end{center}

\begin{center}
\begin{tabular}{ccccccccccc}
\toprule
$i$ & 11 & 12 & 13 & 14 & 15 & 16 & 17 & 18 & 19 & 20 \\
\midrule
$x_i$ & 2.0 & 2.1 & 2.2 & 2.3 & 2.4 & 2.5 & 2.6 & 2.7 & 2.8 & 2.9 \\
$t_i$ & 0.8 & 1.1 & 0.9 & 1.2 & 1.2 & 1.2 & 1.4 & 1.5 & 1.5 & 1.7 \\
\bottomrule
\end{tabular}
\end{center}

\section*{Метод наименьших квадратов}
Ищем такие коэффициенты $w_0,w_1,w_2$, при которых сумма квадратов ошибок
минимальна:
\[
    S(w_0,w_1,w_2)
    =
    \sum_{i=1}^{n}\left(t_i-w_0-w_1x_i-w_2x_i^2\right)^2
    \longrightarrow \min.
\]

В матричной форме задача записывается через матрицу признаков
\[
    X=
    \begin{pmatrix}
        1 & x_1 & x_1^2\\
        1 & x_2 & x_2^2\\
        \vdots & \vdots & \vdots\\
        1 & x_n & x_n^2
    \end{pmatrix},
    \qquad
    w=
    \begin{pmatrix}
        w_0\\
        w_1\\
        w_2
    \end{pmatrix},
    \qquad
    t=
    \begin{pmatrix}
        t_1\\
        t_2\\
        \vdots\\
        t_n
    \end{pmatrix}.
\]
Тогда вектор предсказаний равен $Xw$, а минимизируемая величина принимает вид
\[
    S(w)=\|t-Xw\|^2.
\]
Условие минимума для квадратичной функции приводит к нормальным уравнениям:
\[
    X^TXw=X^Tt.
\]

\section*{Расчет нормальных уравнений}
Для модели $f(x)=w_0+w_1x+w_2x^2$ система нормальных уравнений имеет вид
\[
\begin{pmatrix}
n & \sum x_i & \sum x_i^2\\
\sum x_i & \sum x_i^2 & \sum x_i^3\\
\sum x_i^2 & \sum x_i^3 & \sum x_i^4
\end{pmatrix}
\begin{pmatrix}
w_0\\
w_1\\
w_2
\end{pmatrix}
=
\begin{pmatrix}
\sum t_i\\
\sum x_it_i\\
\sum x_i^2t_i
\end{pmatrix}.
\]

По исходной выборке получаем следующие суммы:
\[
\begin{aligned}
    n &= 20, &
    \sum x_i &= 39.0000, &
    \sum x_i^2 &= 82.7000,\\
    \sum x_i^3 &= 187.2000, &
    \sum x_i^4 &= 444.8666, &
    \sum t_i &= 26.7000,\\
    \sum x_it_i &= 51.0300, &
    \sum x_i^2t_i &= 108.0670. &
\end{aligned}
\]

Следовательно,
\[
\begin{pmatrix}
20 & 39.0000 & 82.7000\\
39.0000 & 82.7000 & 187.2000\\
82.7000 & 187.2000 & 444.8666
\end{pmatrix}
\begin{pmatrix}
w_0\\
w_1\\
w_2
\end{pmatrix}
=
\begin{pmatrix}
26.7000\\
51.0300\\
108.0670
\end{pmatrix}.
\]

Решая данную систему, получаем:
\[
\boxed{
    w_0=4.996625,\qquad
    w_1=-3.929904,\qquad
    w_2=0.967760
}.
\]
Итоговая регрессионная зависимость:
\[
\boxed{
    f(x)=4.996625-3.929904x+0.967760x^2
}.
\]

\section*{Проверка по точкам выборки}
\begin{center}
\small
\begin{tabular}{ccccc}
\toprule
$i$ & $x_i$ & $t_i$ & $f(x_i)$ & $t_i-f(x_i)$ \\
\midrule
1  & 1.0 & 2.0 & 2.034481 & -0.034481 \\
2  & 1.1 & 1.9 & 1.844720 & 0.055280 \\
3  & 1.2 & 1.7 & 1.674314 & 0.025686 \\
4  & 1.3 & 1.6 & 1.523264 & 0.076736 \\
5  & 1.4 & 1.3 & 1.391569 & -0.091569 \\
6  & 1.5 & 1.1 & 1.279229 & -0.179229 \\
7  & 1.6 & 1.3 & 1.186244 & 0.113756 \\
8  & 1.7 & 1.2 & 1.112614 & 0.087386 \\
9  & 1.8 & 1.1 & 1.058340 & 0.041660 \\
10 & 1.9 & 1.0 & 1.023421 & -0.023421 \\
11 & 2.0 & 0.8 & 1.007857 & -0.207857 \\
12 & 2.1 & 1.1 & 1.011648 & 0.088352 \\
13 & 2.2 & 0.9 & 1.034795 & -0.134795 \\
14 & 2.3 & 1.2 & 1.077297 & 0.122703 \\
15 & 2.4 & 1.2 & 1.139154 & 0.060846 \\
16 & 2.5 & 1.2 & 1.220366 & -0.020366 \\
17 & 2.6 & 1.4 & 1.320933 & 0.079067 \\
18 & 2.7 & 1.5 & 1.440856 & 0.059144 \\
19 & 2.8 & 1.5 & 1.580133 & -0.080133 \\
20 & 2.9 & 1.7 & 1.738766 & -0.038766 \\
\bottomrule
\end{tabular}
\normalsize
\end{center}

\section*{Вычисление RMSE}
Сумма квадратов остатков равна
\[
    \sum_{i=1}^{20}\left(t_i-f(x_i)\right)^2 = 0.180189.
\]
Тогда средний квадрат ошибки:
\[
    MSE=\frac{1}{20}\cdot 0.180189=0.009009.
\]
Отсюда
\[
\boxed{
    RMSE=\sqrt{0.009009}=0.094918
}.
\]
Полученное значение RMSE меньше $0.1$, то есть типичное отклонение модели от
наблюдаемых значений составляет примерно $0.095$ единицы по оси $t$.

\section*{Дополнительная оценка качества}
Для интерпретации результата полезно также вычислить общую сумму квадратов
отклонений наблюдений от среднего значения:
\[
    SST=\sum_{i=1}^{20}(t_i-\overline{t})^2,
    \qquad
    \overline{t}=\frac{1}{20}\sum_{i=1}^{20}t_i=1.335.
\]
По данным выборки
\[
    SST=1.985500.
\]
Коэффициент детерминации равен
\[
    R^2=1-\frac{SSE}{SST}
       =1-\frac{0.180189}{1.985500}
       =0.909248.
\]
Следовательно, построенная квадратичная модель объясняет примерно $90.9\%$
разброса значений $t_i$ относительно их среднего значения.

Так как $w_2>0$, парабола направлена ветвями вверх. Ее вершина находится в
точке
\[
    x_*=-\frac{w_1}{2w_2}=2.030412,
    \qquad
    f(x_*)=1.006962.
\]
Это согласуется с таблицей исходных данных: минимальные наблюдаемые значения
расположены около $x=2.0$.

\section*{График регрессии}
\begin{figure}[H]
\centering
\begin{tikzpicture}
\begin{axis}[
    width=0.9\textwidth,
    height=0.55\textwidth,
    grid=both,
    xlabel={$x$},
    ylabel={$t$},
    xmin=0.95,
    xmax=2.95,
    ymin=0.7,
    ymax=2.15,
    legend pos=north east,
    tick label style={font=\small},
    label style={font=\small},
    legend style={font=\small},
]
\addplot[
    only marks,
    mark=*,
    mark size=2pt,
    color=blue
] coordinates {
    (1.0,2.0)
    (1.1,1.9)
    (1.2,1.7)
    (1.3,1.6)
    (1.4,1.3)
    (1.5,1.1)
    (1.6,1.3)
    (1.7,1.2)
    (1.8,1.1)
    (1.9,1.0)
    (2.0,0.8)
    (2.1,1.1)
    (2.2,0.9)
    (2.3,1.2)
    (2.4,1.2)
    (2.5,1.2)
    (2.6,1.4)
    (2.7,1.5)
    (2.8,1.5)
    (2.9,1.7)
};
\addlegendentry{Выборка $(x_i,t_i)$}

\addplot[
    domain=1.0:2.9,
    samples=200,
    thick,
    color=red
] {4.996624515835 - 3.929904306220*x + 0.967760309866*x^2};
\addlegendentry{$f(x)=4.996625-3.929904x+0.967760x^2$}
\end{axis}
\end{tikzpicture}
\caption{Исходные точки и найденная квадратичная регрессия.}
\end{figure}

\section*{График остатков}
Остатки $e_i=t_i-f(x_i)$ показывают, насколько сильно каждое наблюдение
отклоняется от найденной регрессионной кривой. Для хорошей аппроксимации они
должны быть небольшими и не образовывать выраженной систематической структуры.

\begin{figure}[H]
\centering
\begin{tikzpicture}
\begin{axis}[
    width=0.9\textwidth,
    height=0.45\textwidth,
    grid=both,
    xlabel={$x$},
    ylabel={$e_i=t_i-f(x_i)$},
    xmin=0.95,
    xmax=2.95,
    ymin=-0.24,
    ymax=0.15,
    ybar,
    bar width=4pt,
    tick label style={font=\small},
    label style={font=\small},
]
\addplot[
    fill=blue!45,
    draw=blue!70!black
] coordinates {
    (1.0,-0.034481)
    (1.1,0.055280)
    (1.2,0.025686)
    (1.3,0.076736)
    (1.4,-0.091569)
    (1.5,-0.179229)
    (1.6,0.113756)
    (1.7,0.087386)
    (1.8,0.041660)
    (1.9,-0.023421)
    (2.0,-0.207857)
    (2.1,0.088352)
    (2.2,-0.134795)
    (2.3,0.122703)
    (2.4,0.060846)
    (2.5,-0.020366)
    (2.6,0.079067)
    (2.7,0.059144)
    (2.8,-0.080133)
    (2.9,-0.038766)
};
\addplot[
    domain=1.0:2.9,
    samples=2,
    thick,
    color=red
] {0};
\end{axis}
\end{tikzpicture}
\caption{Остатки регрессии для каждой точки выборки.}
\end{figure}

\section*{Вывод}
Методом наименьших квадратов была построена квадратичная регрессионная модель
\[
    f(x)=4.996625-3.929904x+0.967760x^2.
\]
Она отражает U-образную зависимость в данных: значения $t$ сначала убывают при
росте $x$, достигают минимума около середины рассматриваемого интервала, а затем
начинают возрастать. Среднеквадратическое отклонение модели от выборки равно
\[
    RMSE=0.094918.
\]
Коэффициент детерминации $R^2=0.909248$ показывает, что модель объясняет
большую часть изменчивости исходных наблюдений. График остатков подтверждает,
что ошибки аппроксимации невелики: максимальные по модулю отклонения находятся
около $0.2$, а типичная ошибка составляет около $0.095$.
Таким образом, полученная регрессия достаточно точно описывает заданные точки.

\end{document}
